\chapter{Introduction}\label{ch:introduction}

Exacta provides metrics and differentiable losses for two- and three-dimensional image
segmentation, with a focus on medical images. Overlap scores (Dice, IoU, sensitivity and
others), surface-distance scores (Hausdorff distance, HD95, average symmetric surface distance
and normalized surface Dice), physical areas and volumes, probability and calibration scores,
lesion and instance evaluation, dataset-level classification and regression metrics, and the
usual training objectives (Dice, cross-entropy, focal, Tversky, boundary and topology losses)
share one interface. CPU work runs in C++ kernels with OpenMP; supported CUDA work runs in
Triton kernels or in native PyTorch.

\section{Why conventions are arguments}\label{introduction:why}

Two tools can report different numbers for a score that has the same name. HD95 may be the
maximum of the two directional 95th percentiles or one percentile over both directions
together; ASSD may weight every surface voxel equally or by surface area; NSD may count a
distance that equals the tolerance, or not, depending on how the last bit of the distance
rounds. Exacta makes each of these choices an argument with a documented default, so that a
score is defined by its name and its arguments. Chapter~\ref{ch:conventions} states every
convention, Chapter~\ref{ch:tools} gives the arguments that reproduce MONAI, MedPy, DeepMind's
surface-distance and MeshMetrics, and the \emph{reporting signature}
(Section~\ref{metrics:reporting-signatures-since-0.1.0}) writes the full definition of a score in one string to
print next to the number.

\section{What is guaranteed}\label{introduction:guarantees}

The accuracy claims differ between components. The table separates what is argued or checked
exactly from what holds to floating-point rounding. The last column names the test files that
check each claim; Chapter~\ref{ch:tests} says how to run them.

\begin{longtable}{@{}
  >{\raggedright\arraybackslash}p{0.22\linewidth}
  >{\raggedright\arraybackslash}p{0.48\linewidth}
  >{\raggedright\arraybackslash}p{0.24\linewidth}@{}}
\toprule
Component & Guarantee & Checked by \\
\midrule
\endhead
Squared Euclidean distance transform &
The result is the exact floating-point minimum, over the feature voxels, of the sum over the axes,
in axis order, of $\mathrm{fl}\big(\mathrm{fl}((q_i-p_i)h_i)^2\big)$ with every product and sum
rounded once. The C++, Triton and PyTorch implementations return the same bits. &
\texttt{test\_edt.py}: brute-force oracle, property test, lines found from the exactness argument \\
Hausdorff distance, HD95 &
Exact for the discrete definition (the distances between boundary voxels, surface elements or
triangle centres), not for the continuous shape that the mask samples. &
\texttt{test\_surface.py}, \texttt{test\_surfels.py}, \texttt{test\_mesh.py} \\
ASSD, weighted NSD &
Sums over surface elements are correctly rounded (the exact sum rounded once), so the bits do
not depend on element order, thread count or device. &
\texttt{test\_sums.py} \\
Mesh distances &
Distances from triangle centres to the other mesh are minima of binary64 point--triangle
distances over all elements. This is tested against brute force, not proved. &
\texttt{test\_mesh.py} \\
Overlap and volume scores &
Computed from exact integer counts of true positives, false positives and false negatives. &
\texttt{test\_overlap.py}, \texttt{test\_volume.py} \\
Losses &
Precise to the rounding of the stated formula in the stated precision. Their bits can differ
between PyTorch builds and between CPU and GPU. &
\texttt{test\_losses.py}, \texttt{test\_composed\_losses.py} and others \\
\bottomrule
\end{longtable}

An exact score is a property of the discrete definition. Exacta does not claim that a metric
computed on a voxel grid equals the metric of the continuous shapes that the masks represent,
and it does not claim that any loss improves a model. A convergence of tools on one number
needs the same definition and an exact computation of it; Exacta supplies the second and makes
the first explicit.

\section{How to read this manual}\label{introduction:reading}

Chapters~\ref{ch:installation} and~\ref{ch:quickstart} install the library and run a first
example. Chapter~\ref{ch:conventions} defines the surface scores and the policies for empty
masks and ties. Chapters~\ref{ch:metrics} and~\ref{ch:losses} describe the metrics and the
losses; Chapter~\ref{ch:tools} lists the arguments that reproduce other tools, and
Chapter~\ref{ch:performance} reports measured speed. Chapter~\ref{ch:other} lists the modules
that this manual does not describe in detail. Chapters~\ref{ch:driver} and~\ref{ch:tests}
describe the example driver and the test subset for referees, and Chapter~\ref{ch:citing}
says how to cite the software.

Code, argument names and file names are set in \texttt{typewriter}. Spatial arrays have shape
$(B,\ast\text{spatial})$ for class indices and $(B,C,\ast\text{spatial})$ for channels, with $B$
cases and $C$ classes. $P$ is the prediction and $T$ the target. Spacings are in the axis
order of the array, in the unit of the reported distances (millimetres in the examples).
